Nullniveau der Lageenergie: der Landepunkt. Zustand 1: am Start (nur Lageenergie). Zustand 2: beim Absprung (Lageenergie und kinetische Energie). Zustand 3: bei der Landung (nur kinetische Energie). Der Absprung ist nur ein Zwischenzustand: Nur die Höhen zählen, nicht die Richtung ihres Flugs.

\begin{abcliste}
\abc
Sei $h_1$ die Höhe des Anlaufs und $v_1$ die Geschwindigkeit beim Absprung. Zwischen den Zuständen 1 und 2 sinkt sie um $h_1$:
\begin{align}
 m\,g\,h_1 &= \frac{1}{2}\,m\,v_1^2 \\
 v_1 &= \vaF \\
   &= \sqrt{2\cdot\g\cdot\hA} \\
   &= \va \approx \result{\vaP{0}{2}}
\end{align}
Das sind etwa \vakmhP{0}{3}.

\abc
Sei $h_2$ die Höhe des Schanzentischs über dem Landepunkt und $v_2$ die Landegeschwindigkeit. Zwischen den Zuständen 1 und 3 sinkt sie um $h_1 + h_2$:
\begin{align}
 m\,g\,(h_1+h_2) &= \frac{1}{2}\,m\,v_2^2 \\
 v_2 &= \vbF \\
   &= \sqrt{2\cdot\g\cdot(\hA+\hL)} \\
   &= \vb \approx \result{\vbP{0}{2}}
\end{align}
\end{abcliste}
